\documentclass[../../../include/open-logic-section]{subfiles}

\begin{document}

\olfileid[id]{sth}{z}{sep}
\olsection{Pemisahan}

Kita mulai dengan sebuah prinsip untuk menggantikan Komprehensi Naif:

\begin{axiom}[Skema Pemisahan] Untuk setiap formula $\phi(x)$, berikut ini
merupakan suatu aksioma: untuk setiap $A$, himpunan
$\Setabs{x \in A}{\phi(x)}$ ada.
\end{axiom}

Perhatikan bahwa ini bukan satu aksioma tunggal, melainkan suatu
\emph{skema} aksioma. Terdapat \emph{tak berhingga banyaknya} aksioma
Pemisahan; satu untuk setiap formula $\phi(x)$. Skema ini juga dapat (dan
biasanya memang) dituliskan sebagai berikut:

\begin{defish}
Untuk setiap formula $\phi(x)$ yang tidak memuat ``$S$'', berikut ini
merupakan suatu aksioma:
\[
	\forall A \exists S \forall x(x \in S \liff (\phi(x) \land x \in A)).
\]
\end{defish}

Sesuai dengan konvensi yang dicatat pada awal \olref[sth][][]{part},
formula-formula~$\phi$ dalam aksioma Pemisahan dapat mempunyai
parameter.\footnote{Untuk penjelasan tentang maknanya, lihat pembahasan tepat
setelah \olref[sfr][infinite][induction]{natinductionschema}.}

Pemisahan langsung dibenarkan oleh konsepsi kumulatif-iteratif tentang
himpunan yang telah kita uraikan. Untuk melihat alasannya, misalkan $A$
suatu himpunan. Jadi, $A$ dibentuk pada suatu tahap~$S$ (menurut
\stageshier). Karena $A$ dibentuk pada tahap~$S$, semua anggota $A$ dibentuk
sebelum tahap $S$ (menurut \stagesacc). Sekarang, secara khusus,
pertimbangkan semua himpunan yang merupakan anggota $A$ dan juga memenuhi
$\phi$; jelas semua himpunan ini pun dibentuk sebelum tahap~$S$. Jadi, semua
himpunan itu juga dibentuk menjadi himpunan
$\Setabs{x \in A}{\phi(x)}$ pada tahap~$S$ (menurut \stagesacc).

Tidak seperti Komprehensi Naif, prinsip ini menghindari Paradoks Russell.
Kita tidak dapat begitu saja menegaskan keberadaan himpunan
$\Setabs{x}{x \notin x}$. Sebaliknya, dengan \emph{diberikannya} suatu
himpunan~$A$, kita dapat menegaskan keberadaan himpunan
$R_A = \Setabs{x \in A}{x \notin x}$. Namun, semua ini hanya membuktikan
bahwa $R_A \notin R_A$ dan $R_A \notin A$; tidak satu pun di antaranya
sangat mengkhawatirkan.

Namun, Pemisahan mempunyai konsekuensi langsung dan mencolok:

\begin{thm}\ollabel{thm:NoUniversalSet}
Tidak ada himpunan \emph{semesta}; yakni, $\Setabs{x}{x = x}$ tidak ada.
\end{thm}

\begin{proof}
Untuk reductio, andaikan $V$ merupakan himpunan semesta. Maka, menurut
Pemisahan, $R = \Setabs{x \in V}{x \notin x} =
\Setabs{x}{x \notin x}$ ada, bertentangan dengan Paradoks Russell.
\end{proof}

Ketiadaan himpunan semesta---bahkan, sifat terbuka hierarki himpunan---adalah
salah satu gagasan paling mendasar di balik konsepsi kumulatif-iteratif. Jadi,
patut diperlihatkan bahwa secara intuitif kita dapat mencapainya melalui
jalur lain. Himpunan semesta harus menjadi !!a{element} dari dirinya sendiri.
Namun, menurut konsepsi kumulatif-iteratif kita, setiap himpunan muncul (untuk
pertama kalinya) dalam hierarki pada tahap pertama tepat sesudah semua
!!{element}s-nya. Akan tetapi, ini menyiratkan bahwa \emph{tidak ada}
himpunan yang menjadi anggota dirinya sendiri. Sebab, setiap himpunan yang
menjadi anggota dirinya sendiri harus muncul untuk pertama kalinya tepat
sesudah tahap tempat himpunan itu pertama kali muncul, dan hal itu mustahil.
(Dalam \olref[sth][spine][rank]{defnsetrank} kita akan melihat cara membuat
penjelasan ini lebih ketat dengan menggunakan gagasan ``peringkat'' suatu
himpunan. Namun, untuk melakukannya, kita perlu terlebih dahulu memiliki
beberapa aksioma tambahan.)

Berikut beberapa konsekuensi lain dari Pemisahan dan Ekstensionalitas.

\begin{prop}\ollabel{prop:emptyexists}
Jika ada suatu himpunan, maka $\emptyset$ ada.
\end{prop}

\begin{proof}
Jika $A$ suatu himpunan, $\emptyset = \Setabs{x \in A}{x \neq x}$ ada
menurut Pemisahan.
\end{proof}

\begin{prop}
$A \setminus B$ ada untuk setiap himpunan $A$ dan $B$.
\end{prop}

\begin{proof}
$A \setminus B = \Setabs{x \in A}{x \notin B}$ ada menurut Pemisahan.
\end{proof}

Ternyata, irisan yang (hampir) sembarang juga ada:

\begin{prop}\ollabel{prop:intersectionsexist}
Jika $A \neq \emptyset$, maka
$\bigcap A = \Setabs{x}{(\forall y \in A)x \in y}$ ada.
\end{prop}

\begin{proof}
Misalkan $A \neq \emptyset$, sehingga terdapat suatu $c \in A$. Maka
$\bigcap A = \Setabs{x}{(\forall y \in A)x \in y} =
\Setabs{x \in c}{(\forall y \in A)x \in y}$, yang ada menurut Pemisahan.
\end{proof}

Namun, perhatikan syarat $A \neq \emptyset$; sebab $\bigcap \emptyset$ akan
menjadi himpunan semesta secara hampa, bertentangan dengan
\olref{thm:NoUniversalSet}.

\end{document}
